\documentclass{fiche}

\surtitre{Fiche 1 · Nombres et calculs}
\titrefiche{Nombres décimaux et calculs}
\accroche{Tout le reste du programme repose là-dessus. Un calcul mal posé fait rater
un problème de géométrie ou de proportionnalité même quand le raisonnement est juste.}
\niveaux{6\ieme{} · 5\ieme{}}
\priorite{3}
\pourquoi{Outil de tous les autres chapitres : aucun exercice de collège ne s'en passe.}
\duree{25 min}
\domaine{Nombres, calcul et résolution de problèmes}

\begin{document}
\entetefiche

\begin{automatismes}[2]
  \auto[6\ieme]{$\dfrac{1}{10}=0{,}1$ \quad $\dfrac{1}{100}=0{,}01$ \quad $\dfrac{1}{1000}=0{,}001$}
  \auto[6\ieme]{$1=\dfrac{10}{10}=\dfrac{100}{100}=\dfrac{1000}{1000}$}
  \auto[6\ieme]{Multiplier ou diviser un décimal par $10$, $100$, $1000$}
  \auto[5\ieme]{$0{,}6\times 7=4{,}2$ \quad $40\times 0{,}03=1{,}2$}
  \auto[5\ieme]{$2{,}7+1{,}4$ \quad $3{,}4-0{,}8$ de tête}
  \auto[5\ieme]{Factoriser avec les tables : $21=3\times 7$}
  \auto[5\ieme]{Connaître les carrés des entiers de $0$ à $12$, et $10^{3}=1000$}
  \auto[5\ieme]{Connaître les critères de divisibilité par $3$ et par $9$}
  \auto[règle]{Priorités : parenthèses \fleche{} puissances \fleche{} $\times\,\div$ \fleche{} $+\,-$}
  \auto[5\ieme]{$k(a+b)=ka+kb$ \quad et \quad $k(a-b)=ka-kb$}
\end{automatismes}

\section{Lire et écrire un nombre décimal}

\begin{definition}[Nombre décimal]
Un \motcle{nombre décimal} est un nombre qui peut s'écrire sous la forme d'une
\motcle{fraction décimale}, c'est-à-dire avec un dénominateur $10$, $100$, $1000$…
\[
4{,}107 \;=\; \frac{4107}{1000} \;=\; 4+\frac{1}{10}+\frac{0}{100}+\frac{7}{1000}
\]
\end{definition}

\begin{center}
\begin{tikzpicture}[x=13mm,y=7mm]
  \foreach \i/\n in {0/centaines,1/dizaines,2/unités}{
    \draw[trait] (\i,0) rectangle (\i+1,1);
    \node[font=\sffamily\fontsize{6.4}{7}\selectfont,text=encredouce] at (\i+0.5,1.32) {\n};}
  \foreach \i/\n in {3/dixièmes,4/centièmes,5/millièmes}{
    \draw[trait] (\i+0.16,0) rectangle (\i+1.16,1);
    \node[font=\sffamily\fontsize{6.4}{7}\selectfont,text=encredouce] at (\i+0.66,1.32) {\n};}
  \node at (2.5,0.5) {\large 4};
  \node[text=prioritetrois,font=\bfseries] at (3.08,0.12) {\large ,};
  \node at (3.66,0.5) {\large 1};  \node at (4.66,0.5) {\large 0};  \node at (5.66,0.5) {\large 7};
  \draw[accent,line width=0.9pt] (0,0) -- (0,1);
  \node[font=\sffamily\fontsize{7}{8}\selectfont,text=accent] at (1.3,-0.55) {partie entière};
  \node[font=\sffamily\fontsize{7}{8}\selectfont,text=accent] at (4.7,-0.55) {partie décimale};
\end{tikzpicture}
\end{center}

\begin{piege}[Le piège le plus fréquent de toute la 6\ieme]
$3{,}7 > 3{,}17$ ! On ne compare \emph{pas} les parties décimales comme des entiers.\\
\textbf{Méthode sûre :} compléter avec des zéros pour avoir autant de décimales
($3{,}70$ contre $3{,}17$), puis comparer chiffre par chiffre.
\end{piege}

\section{Comparer, encadrer, arrondir}

\textbf{Comparer} : d'abord la partie entière ; si elle est égale, on compare les
dixièmes, puis les centièmes, etc.

\begin{tableaufiche}{@{}>{\bfseries}G{26mm} Y Y Y@{}}
\ligneentete \entetecell{Nombre} & \entetecell{À l'unité} & \entetecell{Au dixième} & \entetecell{Au centième}\\
Encadrer $7{,}368$ & $7 < 7{,}368 < 8$ & $7{,}3 < 7{,}368 < 7{,}4$ & $7{,}36 < 7{,}368 < 7{,}37$\\
Arrondir $7{,}368$ & $\approx 7$ & $\approx 7{,}4$ & $\approx 7{,}37$\\
\end{tableaufiche}

\begin{methode}[Arrondir]
On regarde le chiffre \textbf{juste après} le rang demandé. S'il vaut $5$ ou plus,
on augmente de $1$ le chiffre du rang ; sinon on le garde tel quel.
\end{methode}

\textbf{Intercaler} un nombre entre $2{,}3$ et $2{,}4$ : il suffit d'ajouter une
décimale, par exemple $2{,}35$. \emph{Entre deux décimaux, il y en a toujours une infinité.}

\section{Multiplier et diviser par 10, 100, 1000 — et par 0,1 ; 0,01 ; 0,001}

\begin{tableaufiche}{@{}>{\bfseries}G{17mm} Y Y@{}}
\ligneentete \entetecell{Opération} & \entetecell{Ce que ça fait} & \entetecell{Exemple}\\
$\times 10$   & la virgule se décale d'un rang vers la droite & $4{,}56\times 10=45{,}6$\\
$\times 100$  & deux rangs vers la droite                     & $4{,}56\times 100=456$\\
$\div 10$     & un rang vers la gauche                        & $4{,}56\div 10=0{,}456$\\
$\times 0{,}1$   & \textbf{revient à diviser par 10}          & $4{,}56\times 0{,}1=0{,}456$\\
$\times 0{,}01$  & \textbf{revient à diviser par 100}         & $4{,}56\times 0{,}01=0{,}0456$\\
\end{tableaufiche}

\begin{piege}
Multiplier ne rend pas toujours plus grand : multiplier par $0{,}1$ \emph{diminue}.
\end{piege}

\section{Les quatre opérations}

\subsection{Le vocabulaire (exigé en 5\ieme)}

\begin{tableaufiche}{@{}G{16mm} G{24mm} Y@{}}
\ligneentete \entetecell{Opération} & \entetecell{Le résultat} & \entetecell{Les nombres s'appellent}\\
$a+b$ & une \textbf{somme}    & des \textbf{termes}\\
$a-b$ & une \textbf{différence} & des \textbf{termes}\\
$a\times b$ & un \textbf{produit} & des \textbf{facteurs}\\
$a\div b$ & un \textbf{quotient} & dividende et diviseur\\
\end{tableaufiche}

\subsection{Multiplier deux décimaux}

\begin{methode}
\begin{enumerate}[leftmargin=6mm,label=\textbf{\arabic*.}]
  \item On calcule \textbf{sans les virgules}.
  \item On compte le nombre total de décimales des deux facteurs.
  \item On place la virgule pour en avoir autant dans le résultat.
\end{enumerate}
\vspace{0.8mm}
$2{,}4\times 1{,}5$ : \quad $24\times 15=360$, et $1+1=2$ décimales,
donc $2{,}4\times 1{,}5 = 3{,}60 = 3{,}6$.
\end{methode}

\textbf{Contrôle par ordre de grandeur} — $2{,}4\approx 2$ et $1{,}5\approx 1{,}5$ :
le résultat doit être proche de $3$. C'est le cas. Le programme demande
explicitement ce réflexe de vérification.

\subsection{Diviser}

\begin{itemize}
  \item \textbf{Division euclidienne} (entiers) : $17 = 3\times 5 + 2$, avec toujours
        \textbf{reste $<$ diviseur}.
  \item \textbf{Division décimale} : on continue après la virgule en abaissant des zéros.
  \item \niv{5\ieme} \textbf{Diviser par un décimal} : on multiplie \emph{les deux}
        nombres par $10$, $100$… pour rendre le diviseur entier.
        \[\frac{7{,}2}{0{,}4}=\frac{72}{4}=18\]
\end{itemize}

\subsection{\niv{5\ieme} Multiples, diviseurs et critères de divisibilité}

Si $a = b\times c$, alors $a$ est un \motcle{multiple} de $b$ et de $c$, et $b$ et $c$
sont des \motcle{diviseurs} de $a$.

\begin{tableaufiche}{@{}>{\bfseries}G{22mm} Y >{\itshape}Y@{}}
\ligneentete \entetecell{Divisible par} & \entetecell{Critère} & \entetecell{Exemple}\\
2  & se termine par $0,2,4,6,8$ & $4536$ : oui\\
3  & la \textbf{somme des chiffres} est divisible par 3 & $4{+}5{+}3{+}6=18$ : oui\\
5  & se termine par $0$ ou $5$ & $4536$ : non\\
9  & la \textbf{somme des chiffres} est divisible par 9 & $18$ est divisible par 9 : oui\\
10 & se termine par $0$ & $4536$ : non\\
\end{tableaufiche}

\section{\niv{5\ieme} Priorités opératoires}

\begin{retenir}[L'ordre, dans tous les cas]
\begin{enumerate}[leftmargin=6mm,label=\textbf{\arabic*.}]
  \item les \textbf{parenthèses}, en commençant par les plus intérieures ;
  \item les \textbf{puissances} (carrés, cubes) ;
  \item les \textbf{multiplications et divisions}, de gauche à droite ;
  \item les \textbf{additions et soustractions}, de gauche à droite.
\end{enumerate}
\end{retenir}

\[
7 + 3\times(12-4)\div 2 \;=\; 7+3\times 8\div 2 \;=\; 7+24\div 2 \;=\; 7+12 \;=\; 19
\]

\begin{piege}
$12\div 4\times 3$ vaut $9$, et non $1$ : à priorité égale, on va \textbf{de gauche à droite}.
\end{piege}

\textbf{Traduire un programme de calcul en une seule expression} est un objectif
d'apprentissage explicite de la 5\ieme :

\begin{quote}
« Choisir un nombre, lui ajouter 2, multiplier par 4, retirer 3 »
\quad devient \quad $(n+2)\times 4 - 3$.
\end{quote}

Ici les parenthèses sont \textbf{indispensables} : sans elles, $n+2\times 4-3$ ne
décrit plus le même programme.

\section{\niv{5\ieme} Distributivité}

\begin{retenir}
\[k\times(a+b)=k\times a+k\times b \qquad\qquad k\times(a-b)=k\times a-k\times b\]
\end{retenir}

C'est d'abord un outil de \textbf{calcul mental} :
\[7\times 102 = 7\times(100+2)=700+14=714 \qquad
  25\times 98 = 25\times(100-2)=2500-50=2450\]

\section{\niv{5\ieme} Carrés et cubes}

$a^{2}=a\times a$ (« $a$ au carré ») \qquad $a^{3}=a\times a\times a$ (« $a$ au cube »).

\par\vspace{0.6mm}\noindent\footnotesize
\begin{tabular}{@{}>{\bfseries}l|*{13}{>{\centering\arraybackslash}p{6.6mm}}@{}}
\rowcolor{accent} \entetecell{$n$} & \entetecell{0} & \entetecell{1} & \entetecell{2} & \entetecell{3} & \entetecell{4} & \entetecell{5} & \entetecell{6} & \entetecell{7} & \entetecell{8} & \entetecell{9} & \entetecell{10} & \entetecell{11} & \entetecell{12}\\
$n^{2}$ & 0 & 1 & 4 & 9 & 16 & 25 & 36 & 49 & 64 & 81 & 100 & 121 & 144\\
\end{tabular}\par\vspace{1.2mm}

\begin{piege}
$5+3^{2}=5+9=14$ : la puissance passe \emph{avant} l'addition.
Alors que $(5+3)^{2}=8^{2}=64$.
\end{piege}

\begin{prolongement}
Le programme suggère d'aborder les \textbf{nombres premiers} en prolongement :
le crible d'Ératosthène, et le fait qu'il en existe une infinité — une des plus
anciennes démonstrations connues, dans les \emph{Éléments} d'Euclide.
\end{prolongement}

\begin{videos}
  \video{lzsYQYvEKss}{Numération et nombres décimaux — 6\ieme}{Les Bons Profs}{3:28}{188 000}
  \video{TJH-fiwAt5s}{Calculs avec des priorités — 5\ieme}{Yvan Monka}{4:54}{354 000}
\end{videos}

\end{document}
